\scanpage{320}\mspage{44}
% archive sheet number: 319

i.e. ceci: si $\unsure{\mathcal{F}|F}$ satisfait aussi la condition

4) $\mathcal{F}\cap W\subset W^{\unsure{u|n}}$ (ou ce qui revient au \={m}, comme $\mathcal{F}$ \deleted{\unsure{satisf}} est stable par chgt de base, que $\mathcal{F}\cap W$ est stable par chgt de base),

alors on en conclut
\[
 \mathcal{F}\subset\mathrm{Fib}_{W}.
\]

Or ceci n'est autre que le lemme de K.S. Brown, sous la forme duale de celle donnée plus haut (3.2.1\unsure{;|,} \unsure{page|p.} 20).

Parmi les six ensembles de flèches envisagés dans (4.4.2), les hypothèses 1) 2) 3) sont satisfaites pour $\mathrm{Parf}_{W}$ et $\mathrm{Parf}_{W_{\unsure{0|\bullet}}}$, à l'exclusion des quatre autres. \deleted{\unsure{Une de}} \deleted{\unsure{Elles}} (les autres satisfont à \deleted{\unsure{toutes}} les conditions\unsure{;|,} à la seule exception de la stabilité par composition). On trouve donc \deleted{\unsure{provisoirement}: $\mathrm{Cat},W,W_0$}:

% slanted
\margin{\uline{NB} Pour $\mathrm{Fib}_{W_0}$, ce sont les conditions 1), 2) qui sont vérifiées, et par contre la \unsure{condition} \deleted{\unsure{2)}} 3) qui fait problème. Notons-la $F(W_0)$. Si elle est satisfaite (i.e. si \unsure{Exist|Ex} de $W_0$-fibrations), alors aussi équivalence des conditions $A(\mathrm{Fib}_{W_0})$, $B(\mathrm{Fib}_{W_0})$, $C(\mathrm{Fib}_{W_0})$.}

\uuline{Corollaire} 4.4.4. Les conditions $A(\mathrm{Parf}_{W_{\deleted{\unsure{0}}}})$, $\deleted{\unsure{A}}B(\mathrm{Parf}_{W_{\unsure{\bullet|0}}})$, $C(\mathrm{Parf}_{W})$ sont équivalentes (nous désignons cette condition sur $W$ par $D$). De même les conditions $A(\mathrm{Parf}_{W_0}),B(\mathrm{Parf}_{W_0}),C(\mathrm{Parf}_{W_0})$ \deleted{\unsure{sont équiv}} relatives
